\section{Approach}

\subsection{Preliminary Work and Requirements Gathering}

To ground EndoChron's design in lived experience, we conducted formative sessions remotely with five participants assigned female sex at birth, three with a confirmed endometriosis diagnosis. Table \ref{tab:participant_demographics} lists some demographic characteristics. 1-2 sessions were conducted per participant at this phase, each lasting an average of $\approx$50 minutes.

Sessions were semi-structured and used a master-apprentice framing\cite{Beyer1997ContextualDesign}. They followed contextual inquiry conventions where participants walked us through their current practices for \circled{1} noticing and recording symptoms day-to-day (three participants used at least one app, one journaled, and one did not engage in any self-tracking), \circled{2} making sense of patterns over time, and \circled{3} preparing for clinical encounters.

\begin{table}[h]
  \centering
  \footnotesize
  \begin{tabular}{|c||c|c|c|c|l|}
    \hline
    \textbf{Participant} & \textbf{Age} & \textbf{Endometriosis Dx} & \textbf{Age at Dx} & \textbf{Uses Self-Tracking App(s)} & \textbf{Other Conditions} \\
    \hline
    \hline
    A                    & 53           &                           & --                 & \checkmark                         & Anxiety, Depression       \\
    \hline
    B                    & 34           & \checkmark                & 26                 & \checkmark                         & Pre-diabetic              \\
    \hline
    C                    & 38           &                           & --                 &                                    &                           \\
    \hline
    D                    & 46           & \checkmark                & ``30s''            &                                    &                           \\
    \hline
    E                    & 39           & \checkmark                & 31                 & \checkmark                         & Hashimoto's, Depression   \\
    \hline
  \end{tabular}
  \caption{Participant demographics, self-tracking app usage, and other conditions during Formative Design.}
  \label{tab:participant_demographics}
\end{table}

All sessions were recorded and transcribed with participant consent. All personally identifying information was deleted from all records immediately after analysis\footnote{Our Human Subjects Research statement may be found in Section~\ref{sec:human-subjects}.}. Thematic Analysis of the transcripts yielded the following, broad, recurring patterns:

\begin{enumerate}
  \item \textbf{Input Friction.} Participants described the act of opening a self-tracking app, navigating to the right view/field, and selecting an option as itself an ``\textit{annoying}'' barrier to journaling. One of the three endometriosis participants reported having abandoned self-tracking apps for this reason\footnote{Primarily. Privacy was an equal concern in her case.}. One of two non-endometriosis participants evinced lower annoyance with this input modality.
  \item \textbf{Mists of Time.} All participants described difficulties with recalling not just symptoms but the surrounding context (food, sleep, life events, for instance) when preparing for a clinical visit. Two of three participants who used at least one self-tracking app reported that the visualizations offered were ``\textit{okay}'' and ``\textit{kinda useful}''. The sole endometriosis participant who used a self-tracking app (not specific to her condition) reported difficulty with trying to connect her app's graphical summaries to her flare-ups.
  \item \textbf{Pre-visit Cognitive and Emotional Load.} Four of five participants, including all three with endometriosis, described a pattern of preparing for clinical visits the day before, trying to construct a coherent narrative (``\textit{hit the key points}'') from memories and/or self-tracking data. ``\textit{You're just really worried about what \underline{you're} going to say and what \underline{they're} going to say.}''
  \item \textbf{Loss of Authorship.} Two of three endometriosis participants described instances where a provider summarized their experience to them in a way that felt ``\textit{incongruent}.'' One of them reported correcting their provider at earlier clinical visits (``\textit{He gets me now.}'')
  \item \textbf{Privacy Concerns.} Four of five participants reported concerns with how their data was gathered and used.
\end{enumerate}

Our findings are consistent with the broader qualitative literature on endometriosis communication\cite{Pichon2021Divided,Pichon2025Betrayal,Mardon2023EndometriosisReview} and directly motivate the design goals enumerated in Section~\ref{sec:design-goals}.

\subsection{Flow Model of Endometriosis Communication}

Following Holtzblatt and Beyer's Contextual Design methodology\cite{Beyer1997ContextualDesign}, we modeled the work of communicating about endometriosis as a Flow Model with three principal \textit{entities}, the \textbf{patient}, the \textbf{device}, and the \textbf{provider}, and three \textit{activities} through which health representations propagate: \textbf{capture}, \textbf{reflection / sensemaking}, and \textbf{clinical communication}. While not a Flow Diagram, the Architecture in Figure \ref{fig:architecture} labels all entities and activities.

With current self-tracking applications, and using a simple Flow Model, several gaps may be noted:

\begin{itemize}
  \item Between Embodied Experience $\rightarrow$ Capture, structured-input modalities introduce under/mistranslation.
  \item Between Capture $\rightarrow$ Reflection, existing apps offer limited sensemaking.
  \item Between Reflection $\rightarrow$ Clinical Communication, no artifacts exist in a form the patient owns and curates.
\end{itemize}

\subsection{Theoretical Grounding}

The design of EndoChron is grounded in Distributed Cognition, which maintains that cognitive tasks are not accomplished by an individual mind in isolation but by a socio-technical system in which representations propagate through representational media\cite{furniss2019sociotechnical}.

The clinical visit for an endometriosis patient fits with this theory: months of embodied experience must be `moved' from the patient's body and memory, through some representational medium, into a form the provider can perceive and act on. The gaps we observed in formative interviews occured at the interfaces between the body $\rightarrow$ form, form $\rightarrow$ clinician, clinician's summary $\rightarrow$ patient's confirmation.

EndoChron is designed as a chain of representational media (speech $\rightarrow$ transcript $\rightarrow$ multi-resolution timeline $\rightarrow$ body-map summary) intended to reduce the cognitive load on the patient at the moment of the visit \textit{and} to preserve patient authorship at each transformation. It aims to externalize clinically relevant experiences into persistent visual and linguistic representations, reducing reliance on memory-intensive tasks.

A secondary theoretical grounding is Coiera's Communication-Conversation Model\cite{Coiera2000ConversationModel}, where the patient and provider need to establish \textit{common ground} for clinical work to proceed to the satisfaction of both patient and provider. EndoChron's pre-visit artifact, the Body Map, is intended to be solid ground (established before the visit) that permits limited visit time to be allocated to the \textit{shifting-ground} work for shared decision making.

\subsection{Design Goals}
\label{sec:design-goals}

Four design goals follow from the formative interviews and theoretical commitments described:

\begin{enumerate}[label=\textbf{DG\arabic*.},leftmargin=*,labelsep=0.5em]
  \item \textbf{Eliminate input friction by accepting unstructured speech as the primary input modality.} The patient's account of their experience should be the input to the system, not a translation of that account into checkboxes. Audio is transcribed on-device and discarded. The patient never sees a structured form during capture.
  \item \textbf{Preserve patient authorship across every downstream representation.} The patient must be allowed to inspect, correct, or reject what the system has inferred at every step in the transformation. The system must be positioned as a ``sounding board,'' not an authority\cite{Pichon2025Betrayal}
  \item \textbf{Support sensemaking at the temporal resolution that matches lived experience.} The system must support navigation from the moment level through days, weeks, months, and years. This addresses both the day-to-day resolution at which patients live their condition and the visit-to-visit resolution at which thier providers operate\cite{Pichon2021Divided}.
  \item \textbf{Produce a visit-preparation artifact that the patient controls.} The system must transform longitudinal speech-derived data into a clinically usable representation, that may span arbitrary time periods, that the patient can curate, edit, and annotate with what to disclose at their visit/encounter.
\end{enumerate}

Table~\ref{tab:needs-goals-features} maps each goal to the user need it addresses and to the EndoChron features that implement it. The full RFC-2119 requirements spec may be found in Appendix~\ref{sec:requirements}.

\begin{table}[H]
  \centering
  \footnotesize
  \renewcommand{\arraystretch}{1.5}
  \begin{tabular}{p{0.33\textwidth}|p{0.3\textwidth}:p{0.33\textwidth}}
    \hline
    \textbf{User Need / Gap}                                                                                                                                             & \textbf{Design Goal}                                                                        & \textbf{EndoChron Features}                                                                                                                                                                                                                                                                                                                      \\
    \hline
    \hline
    Patients abandon self-tracking because structured forms create friction and mistranslate embodied experience (Interviews; \cite{McKillop2019Thesis}).                & \textbf{DG1.} Accept unstructured speech as the primary input modality.                     & \textbf{Recording UI:} Voice recording interface, on-device transcription, no structured input imposed \textcolor{phendopink}{(Capture~1--3)}.                                                                                                                                                                                                   \\
    \hline
    Patients distrust algorithmic representations of their experience after experiences of dismissal and misalignment\cite{Pichon2025Betrayal,Ballard2006Endometriosis}. & \textbf{DG2.} Preserve patient authorship across every representation.                      & Transparent display of extracted Phendo terms \textcolor{phendopink}{(Capture~6)}, patient override of all inferences including day rating \textcolor{phendopink}{(Capture~7)}, inline editing of summaries \textcolor{phendopink}{(Reflective Sensemaking~6)}, patient-curated disclosure in visit prep \textcolor{phendopink}{(Visit Prep~6)}. \\
    \hline
    Patients cannot recall context surrounding flare-ups in the visit \& need sensemaking support at multiple temporal scales (Interviews; \cite{Pichon2021Divided}).    & \textbf{DG3.} Support sensemaking at the temporal resolution that matches lived experience. & \textbf{Timeline UI:} Four temporal resolutions (day, week, month, year), LLM-generated per-week, per-month, per condition summaries \textcolor{phendopink}{(Timeline~1--4; Reflective Sensemaking~1--2)}.                                                                                                                                       \\
    \hline
    Patients spend hours preparing for visits \& need to construct narrative \& evidence and control what is disclosed (Interviews; \cite{Pichon2021Divided}).           & \textbf{DG4.} Produce a visit-preparation artifact the patient controls.                    & \textbf{Visual Summarization UI:} Anterior/posterior body map with tappable/clickable severity-colored zones, four lookback time ranges, narrative summary per range, patient controlled disclosure \textcolor{phendopink}{(Visit Prep~1--6)}.                                                                                                   \\
    \hline
    Patients fear data exposure and have privacy concerns.                                                                                                               & \textbf{Cross-cutting.} Preserve privacy as precondition for trust.                         & Exclusively on-device storage, no telemetry, audio discarded after transcription, deidentified transcripts sent upstream for summaries only \textcolor{phendopink}{(Identity \& Privacy~1--4)}.                                                                                                                                                  \\
    \hline
    Patients risk harm if tracking becomes another source of pressure or failure\cite{McKillop2019Thesis}.                                                               & \textbf{Cross-cutting.} Do No Harm in design of intervention.                               & No gamification or pressure to log consistently \textcolor{phendopink}{(Capture~8)}, no language implying clinical certainty or medical advice \textcolor{phendopink}{(Reflective Sensemaking~3)}.                                                                                                                                               \\
    \hline
  \end{tabular}
  \caption{Mapping from User Needs $\rightarrow$ Design Goals $\rightarrow$ Features. \textcolor{phendopink}{Feature numbers} correspond to the Project Requirements in Appendix~\ref{sec:requirements}.}
  \label{tab:needs-goals-features}
\end{table}

\subsection{Prototype}

\subsubsection{Architecture and Implementation}

Figure~\ref{fig:architecture} shows the overall architecture of our intended solution. An interactive prototype of EndoChron is deployed at \href{https://docker.nikhil.io}{\texttt{docker.nikhil.io}} with the source available on \href{https://github.com/afreeorange/endochron}{GitHub}. The prototype implements all four design goals across three views: \textit{Record}, \textit{Reflect}, and \textit{Prepare}. Bruce Tognazzini's \textit{
  First Principles of Interaction Design}\cite{tognazzini2014interaction} were consulted prior to initial design and functioned as a cohesive set of principles throughout the iterative design process.

\begin{figure}[H]
  \centering
  \makebox[\textwidth][c]{%
    \includegraphics[width=1.2\textwidth]{./media/Architecture.jpg}
  }
  \caption{\small Proposed EndoChron Architecture. Users speak freely about symptoms (1), then curate, synthesize, and reflect on AI-generated timelines and body-map summaries (2) before optionally sharing with clinicians or caregivers (3). On-device transcription and vocabulary mapping preserve privacy. De-identified cloud-based summarization and health-rating models support longitudinal synthesis.}
  \label{fig:architecture}
\end{figure}

\subsubsection{Synthetic Data}

At this time, Synthetic Data exclusively drives the prototype (i.e., no audio is actually captured.) The data is meant to simulate a fictitious patient's account of living with endometriosis over a 15-month period. It was generated using Claude Opus 4.7 using the prompt shown in Appendix \ref{sec:prompt}. It may also be examined \href{https://github.com/afreeorange/endochron/blob/main/src/data/syntheticData.json}{on Github}. We simulate $\approx$440 days of app use with a 65\% engagement rate.

\subsubsection{User-Centered Design}

Using an initial prototype, we engaged in iterative, user-centered design, with multiple `live-coding' feedback sessions using the same pool of participants involved in the requirements gathering phase. The current version of EndoChron incorporates almost all UI and UX critique and feedback received during these sessions. Appendix \ref{fig:iphone-screens} shows some screenshots of the final interactive prototype.
